\section{Conclusion}
\label{sec:conclusion}

We tested whether forgetting in continual instruction tuning follows a power law, hits reasoning tasks 3--5$\times$ harder than pattern-matching tasks, and hits early-learned tasks disproportionately.
In naive sequential LoRA tuning of two small LLMs, none of the three holds as stated.
Forgetting is neither uniform nor well described as a per-task power law: it arrives as a fast, saturating hit at each task switch.
Reasoning tasks are not forgotten 3--5$\times$ faster (about 1$\times$ on \llama, about 2$\times$ with wide uncertainty on \qwen), and early-learned tasks are not forgotten disproportionately on \llama once exposure is controlled.
Which specific task follows which is the strongest visible determinant.

The key takeaway is that, in this setting, the forgetting tax is charged per task switch, almost immediately, and its size depends on the pair.

Several follow-ups would sharpen this picture.
First, we would repeat the onset experiment without warm-up, at several learning rates, with more seeds, and on \qwen, to see whether the saturation time scales with learning rate as a step-size effect would.
Second, we would test whether the per-switch hit is predictable from gradient or representation overlap between the two tasks, as the pair matrix suggests.
Third, we would run 15--20 tasks per category so that category can be tested with tasks as the unit, reinstate the full-fine-tuning ablation, and add a larger model.
Finally, we would simulate averaged curves from fitted per-switch steps to check directly whether aggregation alone reproduces the power-law fit.
